Mechanism geometry and kinematics gave me a language for machines: dimensions constrain motion, and transformations connect joint coordinates to an end effector. Yet a useful robot cannot stop at a correct geometric model. Its cameras must produce interpretable signals, software must exchange coherent states, controllers must convert decisions into stable motion, and learning methods must be evaluated against physical behavior. My work has progressed from analyzing how mechanisms move to asking how perception, software, control, and learning can form one dependable machine.

I built this progression through complementary degrees. I completed a B.Eng. in Mechanical Design, Manufacturing and Automation at Harbin Institute of Technology, Weihai, in June 2025. Its training in mechanism behavior, kinematics, and physical constraints established my foundation in robotic systems. I am now pursuing a second bachelor's degree in Computer Science and Technology at Harbin Institute of Technology, Shenzhen, expected in June 2027. This is a deliberate addition to my mechanical preparation: algorithms, data structures, and software interfaces allow me to construct machines whose capabilities extend beyond fixed geometry. Studying computation alongside mechanics lets me reason from physical design through perception and control to integrated operation.

The Robot Digital-Twin HMI for Weld-Seam Recognition joined those domains in one engineering system. For its LeArm representation, I built a modified-DH model, derived analytical forward and inverse kinematics, and cross-checked them with MATLAB Robotics Toolbox. I then loaded binary STL link models in OpenGL and accumulated the modified-DH transform stack so joint inputs updated link and end-effector poses. This visualizer was a kinematic HMI representation, not a physics simulator. Around it, I integrated C++/Qt controls with Python/PyTorch segmentation inference through QProcess, joining image selection, result display, robot interaction, and digital-twin visualization. The experiments likewise required disciplined interpretation. The original U-Net recorded 81.19\% mIoU and 0.64 seam-class IoU, while the final 939-image configuration recorded 96.80\% mIoU and 0.94 seam-class IoU. Because architecture, pretraining, and dataset size changed across the recorded stages, I do not attribute the differences to any single factor or treat them as a controlled ablation. Moreover, the output remained pixel-level segmentation and prompted a fixed preset pose rather than a trajectory derived from seam geometry. These limits showed that connecting perception, geometry, and motion requires explicit interface validation.

That lesson guided my later integration of a Dobot CR5, an Orbbec Astra2 RGB-D camera, and an electric gripper. Eye-to-hand calibration produced a usable static transform for fixed-view RGB data collection and inference, but I documented that high-precision calibration was unfinished because bugs and excessive reprojection error remained. Rather than interpret a usable transform as complete geometric accuracy or as evidence of broader deployment readiness, I treated calibration limits as constraints on subsequent operation. I combined workspace-bound, start-pose-distance, and per-step motion checks with slow playback before using or replaying data. In the ROS2 real-robot loop, ServoP drove the arm and Modbus controlled the gripper. Non-blocking execution created back-and-forth corrections because new inference observed intermediate states before earlier action chunks finished. I therefore adopted blocking, short-horizon ServoP chunks, accepting lower apparent frequency to retain state-action consistency and trajectory quality. The choice reinforced a mechanical-systems principle: integration quality is governed not only by component performance but also by reference frames, admissible motion, timing, and the execution semantics between them.

My Xbotics Embodied AI Community Internship also tested this multidisciplinary discipline. I migrated Isaac Lab's ANYmal-C navigation task into a MotrixLab NumPy environment, refactoring reset and step flows, command sampling, observations, rewards, and termination. I treated the migration as preservation of behavior-defining interfaces rather than translation of code alone. I corrected MuJoCo heightfield frames, constrained spawn and goal locations, and added rollout and reward/termination regression checks. With no recorded success gain or real-quadruped result, these checks verified implementation consistency rather than performance. The experience strengthened my habit of testing coordinate, state, and environment assumptions before trusting an integrated system.

My immediate goal is an R\&D engineering role in robotics, automation, or intelligent manufacturing, connecting mechanics with perception, control, and dependable software. I want to deepen my ability to model physical systems, evaluate data-driven methods, and validate their behavior on real hardware. After gaining stronger real-system experience, I may pursue doctoral study; it remains an option rather than a predetermined route. I therefore seek program-specific training and applied opportunities that can turn this cross-disciplinary foundation into rigorous machine development.
